\setlength{\absparsep}{18pt}
\begin{resumo}[Abstract]
\begin{otherlanguage*}{english}
Practical learning in Operating Systems (OS), particularly within the Linux ecosystem, is a cornerstone in the education of Information Technology professionals. Traditionally, hands-on lab sessions rely on local Virtual Machines (VMs) or dual-boot configurations. However, this model introduces severe barriers: high consumption of RAM and CPU resources, hardware virtualization incompatibilities (such as disabled VT-x in students' BIOS settings), significant instructional time lost to technical troubleshooting instead of core conceptual teaching, and diminished student engagement with static exercise guides. Furthermore, browser-based simulators running entirely on the client side prove unsuitable for formal summative assessments, as evaluation logic can be trivially inspected or altered via browser developer tools (DevTools). This Undergraduate Thesis presents the design and development of an integrated educational platform for practical learning and assessment in Operating Systems. The proposed solution combines a high-fidelity Linux/POSIX simulation engine built in pure TypeScript with an in-memory Virtual File System (VFS) running client-side with zero latency, paired with a secure, decoupled back-end architecture. The back-end orchestrates academic classrooms, time-controlled examinations enabled by dynamic release tokens during class hours, dedicated scheduling for make-up exams, and a server-side evaluation engine powered by signed VFS state snapshots. The initial technical validation comprises an automated test suite with over 750 passing tests covering POSIX compliance and industry certification scenarios. We conclude that this architecture removes local environment friction, maximizes active instructional time in the classroom, and guarantees assessment integrity for the Operating Systems curriculum at UTFPR Guarapuava Campus.

\textbf{Keywords:} Operating Systems. Linux. Web Simulation. Virtual File System. Computer Science Education. Practical Assessment. UTFPR.
\end{otherlanguage*}
\end{resumo}
